\subsection*{6.2 Analysis}
In this section we analyze the LDC obtained above. We prove
\setcounter{nttheorem}{64}
\begin{ntproposition}[Author Proposition 65]
With the notation of the previous section, $C$ is a $(q,\delta,\varepsilon)$-LDC, where
\[q\leq q_1q_2q_3,\qquad\delta\geq\frac{\delta_2\delta_3}{16},\qquad\varepsilon\leq\varepsilon_1+(\varepsilon_2+\varepsilon_3)n.\]
Furthermore, $C$ has rate $\rho_1\rho_2\rho_3$, where $\rho_1,\rho_2$ are as defined in the building blocks paragraph, and per our convention set above, $\rho_3=\rho_3(D/2)$.
\end{ntproposition}
\paragraph{Remark regarding the distance}
Note that the distance $\delta$ of the resulted code $C$ is independent of $\delta_1$, the poor distance of $C_1$ we set out to amplify. This is the key feature of the AEL distance amplification procedure (which our variant above, of course, maintains). It is yet another instance of a general strategy in pseudo-randomness that combines objects in such a way that the resulted object enjoys the upsides of the different parts and avoid their shortcomings. The Zig-Zag product is another classic example. But, of course, $\delta_1$ has some effect on the resulted code. The effect $\delta_1$ has on the code is via the query complexity. Indeed, as the analysis will show, the smaller $\delta_1$ is (i.e., the weaker the guarantee we have on the distance of $C_1$), the larger $k_2=k_2(\delta_1)$ and $k_3=k_3(\delta_1)$ must be, with a far stronger effect on $k_3$. More quantitatively, roughly speaking, by taking a sufficiently good sampler (e.g., the one that is given by Theorem 15), $k_2\approx\operatorname{poly}\log(1/\delta_1)$ and $k_3\approx\operatorname{poly}(1/\delta)$. This, in turn, effects the query complexities $q_2=q_2(k_2)$ and $q_3=q_3(k_3)$.
\begin{proof}[Author proof of Proposition 65]
That the query complexity is $q\leq q_1q_2q_3$ readily follows by the querying function, where recall that per our convention $q_3=q_3(2D)$. To analyze the rate, recall that $\rho_3$ is a non-decreasing function. Further, our convention dictates that by writing $\rho_3$ without explicitly mentioning the message length, we refer to $\rho_3$ applied with the smallest message length taken by the construction, namely, $\rho_3=\rho_3(D/2)$. Thus,
\[n=\sum_{v\in R}n_v=\sum_{v\in R}\frac{k_v}{\rho_3(k_v)}\leq\frac1{\rho_3}\sum_{v\in R}k_v=\frac{\ell n_2}{\rho_3}=\frac{n_1}{\rho_2\rho_3}.\]
Recall that $k=k_1=\rho_1n_1$ which shows that $\rho=k/n\geq\rho_1\rho_2\rho_3$.
We turn to analyze the distance $\delta$ and error $\varepsilon$. Let $x\in\Sigma^k$ and let $\widetilde{C(x)}\in\Sigma^n$ be such that $\operatorname{dist}(C(x),\widetilde{C(x)})\leq\delta$. Define the set of ``errors'', namely, the disagreements between $C(x)$ and $\widetilde{C(x)}$ by
\[B=\{i\in[n]\mid\widetilde{C(x)}_i\ne C(x)_i\}.\]
By assumption, $\mu(B)\leq\delta$. The error set $B$ induces errors ``backwards'' throughout the construction. We proceed by analyzing these induced errors. Recall that, in the encoding function, we defined for each $v\in[r]$ an element $t^{(v)}=t^{(v)}(x)\in\Sigma^{n_v}$. Partition $\widetilde{C(x)}$ to $r$ substrings $\widetilde{C(x)}=\widetilde t^{(1)}\circ\cdots\circ\widetilde t^{(r)}$, where $\widetilde t^{(v)}$ has length $n_v$, and define the set
\[B_t=\{v\in[r]\mid\operatorname{dist}(t^{(v)},\widetilde t^{(v)})\geq\delta_3\}.\]
Informally, $v\in B_t$ if the adversary has introduced too many errors on the respective block to allow for correct decoding via $D_3$.
\begin{ntclaim}[Author Claim 66] $\mu(B_t)\leq8\delta/\delta_3$.\end{ntclaim}
\begin{proof}
For $v\in R$ let $e_v=\operatorname{dist}(t^{(v)},\widetilde t^{(v)})$. We have that $\sum_{v\in R}e_vn_v\leq\delta n$. On the other hand,
\[\sum_{v\in R}e_vn_v\geq\delta_3\sum_{v\in B_t}n_v\geq\frac{\delta_3D|B_t|}{2}.\]
But, per our assumption that $\rho_3\geq1/2$, and since $k_v\leq2D$ for all $v\in R$,
\[n=\sum_{v\in R}n_v\leq2\sum_{v\in R}k_v\leq4Dr.\]
The proof follows by the above three inequalities.
\end{proof}
For convenience we also denote $B_w=B_t$. Next, we define
\begin{equation}B_z=\{u\in[\ell]\mid|\Gamma(u)\cap B_w|\geq\delta_2n_2\}.\tag{6.1}\end{equation}
\begin{ntclaim}[Author Claim 67] $\mu(B_z)\leq\delta_1$.\end{ntclaim}
\begin{proof}
By Claim 66 and by our assumption on $\delta$,
\[\mu(B_w)\leq\frac{8\delta}{\delta_3}=\frac{\delta_2}{2}.\]
Recall that $G$ is a $(\delta_2/2,\delta_1)$-sampler. Thus, at most $\delta_1$-fraction of the left vertices $u\in[\ell]$ satisfy $\mu(\Gamma(u)\cap B_w)\geq\mu(B_w)+\delta_2/2$. The proof then follows since $\mu(B_w)\leq\delta_2/2$.
\end{proof}
Lastly, define
\[B_y=\left\{a\in[n_1]\;\middle|\;\left\lceil\frac a{k_2}\right\rceil\in B_z\right\}.\]
For $v\in[r]$, $b\in[k_v]$ we define the function $\widetilde w_b^{(v)}:\{0,1\}^{r_3}\to\Sigma$ as follows: on input $s_3\in\{0,1\}^{r_3}$
\[\widetilde w_b^{(v)}(s_3)=D_3(b,\widetilde t^{(v)},s_3).\]
\begin{ntclaim}[Author Claim 68]
There exists a set $E_3\subseteq\{0,1\}^{r_3}$ with $\mu(E_3)\leq\varepsilon_3n$ such that for every $s_3\in\{0,1\}^{r_3}\setminus E_3$, $v\in[r]\setminus B_t$, and $b\in[k_3]$ it holds that $\widetilde w_b^{(v)}(s_3)=w_b^{(v)}$.
\end{ntclaim}
\begin{proof}
Fix $v\in[r]\setminus B_t$. By the definition of $B_t$, one has that $\operatorname{dist}(t^{(v)},\widetilde t^{(v)})\leq\delta_3$. By the encoding function, $t^{(v)}=C_3(w^{(v)})$. Therefore, for every $b\in[k_3]$,
\[\Pr_{s_3\sim U_{r_3}}[D_3(b,\widetilde t^{(v)},s_3)\ne w_b^{(v)}]\leq\varepsilon_3.\]
The proof then follows by taking the union bound over all $v\in[r]\setminus B_t$ and $b\in[k_v]$ as indeed $\sum k_v\leq n$.
\end{proof}
For $(u,j)\in[\ell]\times[n_2]$ we define the function $\widetilde z_j^{(u)}:\{0,1\}^{r_3}\to\Sigma$ as follows. For $s_3\in\{0,1\}^{r_3}$ we have $\widetilde z_j^{(u)}(s_3)=\widetilde w_{j'}^{(v)}(s_3)$, where $(v,j')=\Gamma(u,j)\in[r]\times[k_v]$. Further define the function $\widetilde z^{(u)}:\{0,1\}^{r_3}\to\Sigma^{n_2}$ by
\[\widetilde z^{(u)}(s_3)=\widetilde z_1^{(u)}(s_3)\circ\cdots\circ\widetilde z_{n_2}^{(u)}(s_3).\]
\begin{ntclaim}[Author Claim 69]
For every $u\notin B_z$ and $s_3\in\{0,1\}^{r_3}\setminus E_3$ it holds that
\[\operatorname{dist}(\widetilde z^{(u)}(s_3),z^{(u)})\leq\delta_2.\]
\end{ntclaim}
\begin{proof}
Fix $s_3\in\{0,1\}^{r_3}\setminus E_3$ and consider any $u\in[\ell]\setminus B_z$. By the encoding function, for every $j\in[n_2]$ it holds that $z_j^{(u)}=w_{j'}^{(v)}$, where $(v,j')=\Gamma(u,j)$. As $v\notin B_z$, at most $\delta_2n_2$ of $j\in[n_2]$ satisfy $v\in B_w$. For every other $j$,
\[\widetilde z_j^{(u)}=\widetilde w_{j'}^{(v)}(s_3)=w_{j'}^{(v)}=z_j^{(u)},\]
proving the claim.
\end{proof}
For $u\in[\ell]$, $a\in[k_2]$ we define the function $\widetilde y_a^{(u)}:\{0,1\}^{r_2}\times\{0,1\}^{r_3}\to\Sigma$ as follows. On $(s_2,s_3)\in\{0,1\}^{r_2}\times\{0,1\}^{r_3}$,
\[\widetilde y_a^{(u)}(s_2,s_3)=D_2(a,\widetilde z^{(u)}(s_3),s_2).\]
\begin{ntclaim}[Author Claim 70]
There exists a set $E_2\subseteq\{0,1\}^{r_2}$ with $\mu(E_2)\leq\varepsilon_2n$ such that for every $u\in[\ell]\setminus B_z$, $a\in[k_2]$, and $(s_2,s_3)\in(\{0,1\}^{r_2}\setminus E_2)\times(\{0,1\}^{r_3}\setminus E_3)$ it holds that $\widetilde y_a^{(u)}(s_2,s_3)=y_a^{(u)}$.
\end{ntclaim}
\begin{proof}
Fix $u\in[\ell]\setminus B_z$. By the encoding function $z^{(u)}=C_2(y^{(u)})$. Recall that
\[\widetilde y_a^{(u)}(s_2,s_3)=D_2(a,\widetilde z^{(u)}(s_3),s_2).\]
As $s_3\notin E_3$, $u\notin B_z$, Claim 69 implies $\operatorname{dist}(\widetilde z^{(u)}(s_3),z^{(u)})\leq\delta_2$. Therefore
\[\Pr_{s_2\sim U_{r_2}}[D_2(a,\widetilde z^{(u)}(s_3),s_2)\ne y_a^{(u)}]\leq\varepsilon_2.\]
The proof then follows by taking the union bound over all $a\in[k_2]$ and $u\in[\ell]\setminus B_z$, and noting that $k_2\ell=n_1\leq n$.
\end{proof}
\begin{ntclaim}[Author Claim 71]
For every $(s_2,s_3)\in(\{0,1\}^{r_2}\setminus E_2)\times(\{0,1\}^{r_3}\setminus E_3)$, it holds that
\[\operatorname{dist}(\widetilde y(s_2,s_3),y)\leq\delta_1,\]
where $\widetilde y(s_2,s_3)$ is the concatenation of the $k_2$-length strings $(\widetilde y^{(u)}(s_2,s_3)\mid u\in[\ell])$.
\end{ntclaim}
\begin{proof}
Note that by Claim 70, the projection of the two strings $\widetilde y(s_2,s_3),y$ to a block corresponding to $u\notin B_z$ are in full agreement. The proof then follows by Claim 67.
\end{proof}
We now conclude the proof of Proposition 65. Let $p\in[k]$, by Claim 71, for every $(s_2,s_3)\in(\{0,1\}^{r_2}\setminus E_2)\times(\{0,1\}^{r_3}\setminus E_3)$, we have that $\operatorname{dist}(\widetilde y(s_2,s_3),y)\leq\delta_1$. Since by the encoding function $y=C_1(x)$, it holds
\[\Pr_{s_1\sim U_{r_1}}[D_1(p,\widetilde y(s_2,s_3),s_1)\ne x_p]\leq\varepsilon_1.\]
The proof then follows since $\mu(E_2)\leq\varepsilon_2n$ and $\mu(E_3)\leq\varepsilon_3n$.
\end{proof}
